% ========================== 4. Experimental program ====================
\section{Experimental Program}\label{sec:methods}

\subsection{Overview}\label{sec:methods-overview}

The program follows the hierarchy of Section~\ref{sec:intro}: state
sufficiency (Phase~1), local fluctuation geometry (Phase~2),
target-relevant retrieval (Probe~1), and identifiability of
query-conditioned routing (Probe~2).  All experiments, seeds,
parameters, audit rules and results are frozen and recorded in the
accompanying research archive; this paper reports only archived results.
Statistical procedures are summarized in
Section~\ref{sec:methods-stats}.

\subsection{Phase 1: state sufficiency}\label{sec:methods-phase1}

\textbf{Question.} Is the scalar membrane potential $h_t$ a sufficient
Markov state of TAN?

\textbf{Natural-collision protocol.} $N = 160{,}000$ random histories of
length $T-1 = 9$ were drawn from four stimulus families (Gaussian white
noise, random steps, ramps, sine waves).  For each model separately, all
unordered pairs $(A,B)$ whose subthreshold potentials at $t = T-1$
collide -- $0.05 < h_{T-1} < \theta$ on both sides and
$|h_{T-1}^{(A)} - h_{T-1}^{(B)}| < \delta = 10^{-4}$ -- form the
\textit{natural-collision ensemble}.  Restricting to the subthreshold
band excludes trivial reset coincidences ($h=0$).  At $t = T = 10$ an
identical probe $x^{*} \sim \mathcal{U}(0.5, 3.0)$ is injected into both
histories and the one-step divergence is measured on the continuous
branch (pre-reset potential $u_{T} = \lambda h_{T-1} + A_T$), i.e.\
\begin{equation}
  \Delta\tilde{h}_T = |u_T^{(A)} - u_T^{(B)}|,
  \qquad K = \Delta\tilde{h}_T / |h_{T-1}^{(A)} - h_{T-1}^{(B)}| .
\end{equation}
Post-reset divergence, spike-outcome classes and reset
re-synchronisation (a common double spike at $T$, then identical
continuation inputs) are tabulated separately.  Under the 1D membrane
state-sufficiency hypothesis (where $u_T = \lambda h_{T-1} + f(x_T)$),
divergence contracts with $K = \lambda < 1$ (LIF predicts $K = \lambda$
exactly on every collision pair); large $K \gg 1$ refutes the sufficiency
of the 1D membrane potential $h$ as a complete state descriptor.

\subsection{Phase 2: local fluctuation geometry}\label{sec:methods-phase2}

\textbf{Question.} Does surprise-gated weighting reorganize the local,
event-locked response geometry relative to uniform temporal integration?

\textbf{Protocol.} Sparse episodic stimuli (Poisson-like onset gaps with
burst structure and amplitude diversity; $T = 5000$, burn-in 500, three
fixed seeds) drive B1--B4 and an input-only control C0.  For each model a
response trajectory is embedded in frozen coordinate frames:
$z_C = [h,\mu,S,A]$ for all models (B1's $A$ is its input, a null
convention), and $z_F = [h,\mu,S,A,\alpha_1,\dots,\alpha_5,C]$ for B3/B4
(matched frame).  Coordinates are globally $Z$-scored per seed
(post-burn-in) and per-seed centred before cross-seed pooling; no
whitening is applied.  Events are aligned at episode onset and analysed
at fixed lag $\tau \in [-6, 30]$; the primary response is the
baseline-subtracted state $\Delta z_{e,\tau}$ relative to a five-step
pre-event window, pooled across events (never per-window PR averaging).
At each lag we compute the effective rank of the pooled response
covariance,
\begin{equation}
  d_G(\tau) = \frac{\operatorname{Tr}(C_z)^2}
  {\operatorname{Tr}(C_z^2)},
\end{equation}
its log-trace $\log\operatorname{Tr}(C_z)$ and eigenvalue spectrum, and
the velocity geometry
\begin{equation}
  d_D(\tau) = \frac{\operatorname{Tr}(C_v)^2}
  {\operatorname{Tr}(C_v^2)}, \qquad
  v_{e,\tau} = z_{e,\tau+1} - z_{e,\tau},
\end{equation}
where $C_v$ is the pooled covariance of same-event adjacent-phase
increments.  $d_{PR}$ is reported and interpreted strictly as an
\emph{effective rank of the local fluctuation covariance}; it is never
called an intrinsic or manifold dimension \citep{gaoganguli2015,
mazzucato2016,facco2017,chungabbott2021}.  Zero-variance windows return
NaN (never a fabricated value).  Amplitude controls are computed for
every lag: raw pooled, amplitude-stratified (median split on the
log-amplitude, pooled-within covariance) and per-coordinate
residualisation on log-amplitude.  Quiet-reference samples lie $\geq 12$
steps from any pulse.  Recovery times ($\tau_{1/2}$ on $d_D$,
$\tau_{e\text{-fold}}$ on $\log\operatorname{Tr}$) are estimated on
isolated events with event bootstrap confidence intervals.

\subsection{Probe 1: temporal distractor task}\label{sec:methods-probe1}

\textbf{Question.} Does stored history contribute to target-relevant
retrieval, as opposed to mere storage?

\textbf{Task (frozen v3).} Each trial has five event steps: a target at
fixed step 5 whose amplitude class $c_T$ is low
($\mathcal{U}(0.55,1.05)$) or high ($\mathcal{U}(0.95,1.45)$) with
overlapping supports; either three distractors at steps 6--8 (drawn from
the same class mixture, so a distractor may exceed the target) or
silence; and a query at step 9 that acts as a \emph{retrieval cue} but
does not enter the label.  The label is three-way:
$c_T \in \{\text{lo},\text{hi}\}$ or \emph{absent} (catch, one third of
trials).  The query re-opens the surprise gate so the window (which still
contains the target at step 9) can be re-read into the response.  A
linear three-class softmax readout is trained on response-trajectory
features $F = [u_9,\dots,u_{13},\, y_9,\dots,y_{13}]$; per-step
observables before the query are never visible (no amplitude leakage at
arrival).  Readout protocol, training (clean + distractor mixed,
$n_\text{train} = 4000$), testing ($n_\text{test} = 2000$) and the three
seeds are identical across B1--B4.  Primary endpoint:
$\Delta_{\text{dist}} = \text{acc}(\text{clean}) -
\text{acc}(\text{distractor})$; main contrast B4 vs.\ B3.

\subsection{Probe 2: query-conditioned binding identifiability audit}
\label{sec:methods-probe2}

\textbf{Question.} Does the current query determine which historical
event is selected ($Q \to K \to V$)?

\textbf{Audit design.} Two events $E_A = (K_A, V_A)$, $E_B = (K_B, V_B)$
with key amplitudes $K_A = 1.0$, $K_B = 2.0$ and values
$V \in \{-1,+1\}$ (fair) are encoded into the scalar stream as
$x(E) = K_{\text{amp}} + V\Delta$, $\Delta = 0.25$; the two events
occupy two uniformly random history slots (of four), the remaining slots
carry noise $\mathcal{N}(0, 0.05^2)$, and the query is
$x_t = 3.0$ ($q = A$) or $4.0$ ($q = B$).  The answer is
$y = V_q$.  Identifiability gates: (A) a linear readout on
$z = [x_{t-4},\dots,x_t]$ (B2 shortcut test) must not exceed 0.58
balanced accuracy; (B) the same readout on B3's uniform-integration
response features must not exceed 0.58; (C) for B4, attention must hit
the target history position in $\geq 50\%$ of trials (chance 25\%),
show target/distractor contrast $R = \mathbb{E}[\alpha_{\text{target}}]/
\mathbb{E}[\alpha_{\text{distractor}}] \geq 2$, and produce
query-dependent attention (Jensen--Shannon divergence between
$\alpha\,|\,q{=}A$ and $\alpha\,|\,q{=}B$ distributions with permutation
$p < 0.01$).  Counterfactual audits permute positions, values and
queries, and hold history fixed while changing the query.  All gates,
thresholds and the final status are frozen:
\texttt{TERMINATED / AUDIT\_FAILED} is the archived outcome of this
audit, i.e.\ an identifiability boundary result, not an incomplete
experiment.

\subsection{Statistical and reproducibility procedures}\label{sec:methods-stats}

All experiments use fixed seeds and frozen parameters recorded in the
archive (\texttt{docs/REPRODUCIBILITY\_LEDGER.md}).  Confidence intervals are
event/trial bootstraps (2000 resamples) unless stated.  Calibration
audits (generator sanity, hidden-label checks, shuffled-label readouts,
class-coverage, permutation nulls for shuffle statistics) are executed
before formal runs; violations stop the run.  Probe-1's three-seed data
were re-analysed after an audit-statistics amendment (empirical
permutation null for the shuffle audit; per-cell coverage statuses
instead of run-blocking flags) on the frozen data alone, with the
deterministic reproduction verified against the frozen per-seed numbers
(maximum deviation $4.7\times10^{-4}$, rounding only).  No result in
this paper was obtained by post-hoc tuning.
